\subsection*{3.3. Propiedades de la categoría cofibrada
$\underline\Psi_{L_\bullet}$}

En esta sección, que no es esencial para comprender lo que sigue,
utilizamos libremente las nociones de [4], con el objeto de comparar
los resultados precedentes con los de ese trabajo.

\underline{Teorema} 3.3.1. \underline{La categoría cofibrada aditiva}
$\underline\Psi_{L_\bullet}$ \underline{sobre} $\underline C$
\underline{es exacta a izquierda} ([4, 2.2]), \underline{admite un
objeto cuasimaximal} ([4, 4.4]) \underline{y, para todo objeto}
$E=(X,\alpha)$ \underline{de esta categoría, existe un complejo
típico} $L^E_\bullet$ \underline{para} $E$ ([4, 3.1]).
\underline{Además, el procomplejo típico de}
$\underline\Psi_{L_\bullet}$ ([4, 3.7 y 4.4])
\underline{es canónicamente isomorfo en} $D(\underline C)$
\underline{al complejo truncado} $[L_\bullet]$
\underline{considerado en} 3.1.

Explicitamos ahora el significado de las distintas afirmaciones de
3.3.1.

3.3.2. Decir que la categoría cofibrada
$\underline\Psi_{L_\bullet}$ es \underline{exacta a izquierda}
significa que es aditiva y que además satisface la condición siguiente:

3.3.2.1. Para toda sucesión exacta
\[
 0\longrightarrow A'\xrightarrow{u}A\xrightarrow{v}A''
 \longrightarrow0
\]
en $\underline C$, el funtor natural (con la notación de [4, 1.8])
\[
 \underline\Psi_{L_\bullet}(\sisi{A}{A'})
 \longrightarrow
 \underline\Psi_{L_\bullet}([A\longrightarrow A''])
\]
es una equivalencia de categorías.

Recordemos que, si $\underline\Psi$ es una categoría cofibrada sobre
$\underline C$ y $v:A\longrightarrow A''$ es un morfismo de
$\underline C$, entonces $\underline\Psi([A\longrightarrow A''])$ es
la categoría de los pares $(E,\rho)$ en los que $E$ es un objeto de
$\underline\Psi(A)$ y $\rho$ es una ``trivialización'' del objeto
correspondiente $v_*(E)$ de $\underline\Psi(A'')$; es decir, un
isomorfismo de este con el objeto cero (2.5.1) de
$\underline\Psi(A'')$. La definición del funtor
\[
 \underline\Psi(A')\longrightarrow
 \underline\Psi([A\longrightarrow A''])
\]
de 3.3.2.1 es evidente. Cuando
$\underline\Psi=\underline\Psi_{L_\bullet}$, los objetos de
$\underline\Psi_{L_\bullet}([A\longrightarrow A''])$ son
explícitamente los triples $(X,\alpha,\rho)$, donde $(X,\alpha)$ es
como en (3.1.3) y $\rho:X\longrightarrow A''$ es un morfismo de
$\underline C$ que satisface
\begin{equation}\tag{*}
 \rho i=v,\qquad \rho\alpha=0.
\end{equation}
Si $v:A\longrightarrow A''$ pertenece a una sucesión exacta corta como
en \sisi{3.2.3}{3.3.2.1}, entonces, como es bien sabido, dar un
morfismo que satisfaga la primera relación (*) equivale, por la
exactitud a izquierda de la categoría cofibrada de extensiones de
$L_0$ por objetos variables de $\underline C$, a dar una extensión
$X'$ de $L_0$ por $A'\simeq\Ker v$ que sea una subextensión de $X$:
se toma $X'=\Ker\rho$. La segunda relación (*) dice que $\alpha$ se
factoriza a través de un morfismo
$\alpha':L_1\longrightarrow X'$. Salvo las compatibilidades cuya
verificación es inmediata, esto demuestra la exactitud a izquierda
enunciada en \sisi{3.2.3}{3.3.2.1}.

3.3.3. Recordemos que un objeto $E_0$ de una categoría cofibrada
$\underline\Psi$ sobre $\underline C$ se denomina
\underline{cuasimaximal} si, para todo objeto $E$ de
$\underline\Psi$ situado sobre un objeto $A$ de $\underline C$,
existe un morfismo $E\longrightarrow E'$ en $\underline\Psi$ que
induce un \underline{monomorfismo} $A\longrightarrow A'$ en
$\underline C$, tal que $\sisi{E}{E'}$ esté \underline{dominado por} $E_0$, es
decir, que exista un morfismo $E_0\longrightarrow E'$. Si
$\underline C$ tiene suficientes objetos inyectivos, esto equivale a
decir que todo objeto de $\underline\Psi$ sobre un objeto inyectivo de
$\underline C$ está dominado por $E_0$.

Para $\underline\Psi=\underline\Psi_{L_\bullet}$ existe el siguiente
objeto cuasimaximal canónico. Pongamos
\[
 L'_1=\Coker(d_1:L_2\longrightarrow L_1)
\]
y sea $\bar d_0:L'_1\longrightarrow L_0$ el morfismo inducido por
$d_0:L_1\longrightarrow L_0$. Sea $E_0=(X_0,\alpha_0)$ el objeto de
$\underline\Psi_{L_\bullet}(L'_1)$ para el cual
$X_0=L_0\oplus L'_1$ es la extensión trivial de $L_0$ por $L'_1$ y
$\alpha_0$ es la composición
\[
 L_1\twoheadrightarrow L'_1
 \xrightarrow{(\bar d_0,\,\mathrm{id}_{L'_1})}
 L_0\oplus L'_1.
\]

Para ver que $E_0$ es cuasimaximal, obsérvese que para todo objeto
$(X,\alpha)$ de $\underline\Psi_{L_\bullet}$ se puede encontrar un
morfismo de él a un objeto $(X',\alpha')$, que induce un
\underline{monomorfismo} $A\longrightarrow A'$ en $\underline C$,
tal que la extensión $X'$ de $L_0$ por $A'$ sea
\underline{escindida}. En efecto, basta tomar $A'=X$ y
$(X',\alpha')=i_*(X,\alpha)$, con la notación de (3.1.2). Por otra
parte, es inmediato que todo objeto $(X',\alpha')$ de
$\underline\Psi$ para el cual la extensión $X'$ sea escindida está
dominado por $E_0$.

3.3.4. Sea $E$ un objeto de una categoría cofibrada aditiva
$\underline\Psi$ sobre $\underline C$, situado sobre un objeto $A$ de
$\underline C$, y consideremos el siguiente funtor del objeto variable
$B$ de $\underline C$:
\begin{equation}\tag{3.3.4.1}
 B\longmapsto\Hom(E,\theta_B),
\end{equation}
donde $\theta_B$ denota el objeto trivial de $\underline\Psi(B)$.
Decimos que $E$ \underline{admite un complejo típico} si este funtor
es representable por un objeto $\Omega_E$ de $\underline C$. Como el
funtor en cuestión se aplica canónicamente en el funtor
$B\longmapsto\Hom(A,B)$ representado por $A$, se obtiene un morfismo
canónico
\begin{equation}\tag{3.3.4.2}
 A\longrightarrow\Omega_E,
\end{equation}
y el \underline{complejo de cadenas} correspondiente de longitud uno
$A\longrightarrow\Omega_E$ en $\underline C$ se denomina
\underline{complejo típico de} $E$.

Cuando $\underline\Psi=\underline\Psi_{L_\bullet}$ y
$E=(X,\alpha)$, se comprueba inmediatamente que el funtor (3.3.4.1)
está representado por $\Coker\alpha$, y que el morfismo canónico
(3.3.4.2) es simplemente la composición
\[
 A\xrightarrow{\ \sisi{j}{i}\ }X
 \xrightarrow{\ \mathrm{can}\ }\Coker\alpha.
\]

3.3.5. Si una categoría cofibrada aditiva $\underline\Psi$ sobre
$\underline C$ admite un objeto cuasimaximal (3.3.3), y existen los
complejos típicos de sus objetos, entonces el complejo típico
$T_\bullet$ de un objeto cuasimaximal $E_0$, considerado como objeto de
la \underline{categoría derivada} $D(\underline C)$, es independiente,
salvo isomorfismo canónico, de la elección de $E_0$ ([4, 4.4]). Se
denomina el \underline{complejo típico de la categoría cofibrada
aditiva} $\underline\Psi$. Si, además, $\underline\Psi$ es exacta a
izquierda (3.3.2), entonces $\underline\Psi$ puede reconstruirse a
partir del complejo típico $T_\bullet$, salvo una equivalencia de
categorías fibradas sobre $\underline C$ única salvo isomorfismo único
([4, 6.10]). También hay isomorfismos canónicos, funtoriales en
$A\in\underline C$ ([4, 5.4]),
\begin{equation}\tag{3.3.5.1}
 \Psi^0(A)\xrightarrow{\sim}\Ext^0(T_\bullet,A),
 \qquad
 \Psi^1(A)\xrightarrow{\sim}\Ext^1(T_\bullet,A).
\end{equation}

Para completar la comparación con [4], demostraremos que cuando
$\underline\Psi=\underline\Psi_{L_\bullet}$ su complejo típico es el
truncamiento $[L_\bullet]$. Esto da otra definición de los
isomorfismos (3.1.8) y (3.2.4) en términos de (3.3.5.1), y se comprueba
que es compatible con las definiciones de 3.1 y 3.2.

Para esta afirmación podemos restringirnos claramente al caso en que
$L_\bullet$ ya esté truncado en grado $1$, pues en 3.1 se observó que
el truncamiento no cambia $\underline\Psi_{L_\bullet}$ salvo
isomorfismo de categorías cofibradas. La construcción explícita de
3.3.4 muestra entonces inmediatamente que el complejo típico del
objeto cuasimaximal descrito en 3.3.3 es canónicamente isomorfo a
$L_\bullet$.

\underline{Observación} 3.3.6. Sea
\[
 0\longrightarrow A'\longrightarrow A\longrightarrow A''
 \longrightarrow0
\]
una sucesión exacta en $\underline C$. Da una sucesión exacta de
grupos $\Ext^i$
\[
 0\longrightarrow\Ext^0(L_\bullet,A')
 \longrightarrow\Ext^0(L_\bullet,A)
 \longrightarrow\cdots\longrightarrow\Ext^1(L_\bullet,A)
 \longrightarrow\Ext^1(L_\bullet,A''),
\]
que, por supuesto, podría continuarse a la derecha introduciendo
$\Ext^i$ para $i\geq2$. Utilizando los isomorfismos (3.1.8) y (3.2.4),
obtenemos la sucesión exacta
\begin{multline}\tag{3.3.6.1}
 0\longrightarrow\Psi^0_{L_\bullet}(A')
 \longrightarrow\Psi^0_{L_\bullet}(A)
 \longrightarrow\Psi^0_{L_\bullet}(A'')
 \xrightarrow{\ \sisi{o}{\partial}\ }\Psi^1_{L_\bullet}(A')\\
 \longrightarrow\Psi^1_{L_\bullet}(A)
 \longrightarrow\Psi^1_{L_\bullet}(A'').
\end{multline}
Todas las flechas salvo $\partial$ están definidas simplemente por la
funtorialidad de $\Psi^0_{L_\bullet}$ y $\Psi^1_{L_\bullet}$. Por otra
parte, [4, 2.5.2] define una sucesión exacta (3.3.6.1), para toda
categoría cofibrada exacta a izquierda $\underline\Psi$ sobre
$\underline C$, mediante un morfismo
\begin{equation}\tag{3.3.6.2}
 \partial:\Psi^0(A'')\longrightarrow\Psi^1(A').
\end{equation}
Este asocia a un elemento $u$ del miembro izquierdo la clase, en el
miembro derecho, del objeto de $\underline\Psi(A')$, determinado salvo
isomorfismo canónico, cuya imagen en
$\underline\Psi([A\longrightarrow A''])$ (notación de 3.3.2) es el
par $(X_0,u)$, donde $X_0$ es el objeto trivial de
$\underline\Psi(A)$.

Para $\underline\Psi=\underline\Psi_{L_\bullet}$, el morfismo
(3.3.6.2) definido por esta \underline{construcción ``geométrica'' es
el opuesto del morfismo de borde} $\partial$ \underline{de} (3.3.6.1).
Así, después de las identificaciones anteriores y de un cambio de
signo en $\partial$, esta última sucesión exacta es precisamente la de
[4, 2.5.2]. Esto se deduce de la última afirmación de [4, 5.4], a la
vista de las compatibilidades señaladas en 3.3.5.

\underline{Observación} 3.3.6.3. Aprovechamos la ocasión para aclarar
un punto de [4, n.º 3.11] relativo a los isomorfismos fundamentales
(3.11.3) y (3.11.4), que desempeñan un papel esencial en la
demostración de [4, 5.4]. El diagrama de la página 35 de [4] es
conmutativo, no anticonmutativo, y por tanto no define tal como está un
morfismo
\[
 u:[L^X_1\longrightarrow L^X_0]\longrightarrow K^\bullet[1].
\]
Debe hacerse anticonmutativo sustituyendo $u^1$ por $-u^1$. También se
podría adoptar la convención opuesta y sustituir $u^0$ por $-u^0$;
sin embargo, la convención indicada, en la que $u^1$ se convierte en
$-u^1$, es aquella para la cual (3.11.3) generaliza el isomorfismo
(3.5.3) de [4], como allí se afirma. Con esta aclaración, también debe
corregirse la afirmación de [4, 5.4]: los dos morfismos de borde allí
considerados \underline{difieren en un signo}.

3.3.7. En [4, 6.13] también se explicitó la varianza de la categoría
cofibrada $\underline\Psi_{L_\bullet}$ sobre $\underline C$ al variar
$L_\bullet$, reconstruyendo, en particular, la categoría de funtores
cocartesianos de $\underline\Psi_{L_\bullet}$ a
$\underline\Psi_{L'_\bullet}$ a partir de los complejos de cadenas
$L_\bullet$ y $L'_\bullet$. El conjunto de clases de isomorfismo de
tales funtores cocartesianos está canónicamente en biyección con
\[
 \Hom_{D(\underline C)}(L'_\bullet,L_\bullet).
\]
Solo utilizaremos una pequeña parte de estos resultados, que
desarrollamos en la sección siguiente mediante la descripción sencilla
de $\underline\Psi_{L_\bullet}$ dada en 3.1.
